\documentclass{amsart}
% For dealing with references we use the comment environment
\usepackage{verbatim}
\newenvironment{reference}{\comment}{\endcomment}
%\newenvironment{reference}{}{}
\newenvironment{slogan}{\comment}{\endcomment}
\newenvironment{history}{\comment}{\endcomment}

% For commutative diagrams we use Xy-pic
\usepackage[all]{xy}

% We use 2cell for 2-commutative diagrams.
\xyoption{2cell}
\UseAllTwocells

% We use multicol for the list of chapters between chapters
\usepackage{multicol}

% This is generally recommended for better output
\usepackage{lmodern}
\usepackage[T1]{fontenc}

% For cross-file-references

% Package for hypertext links:
\usepackage{hyperref}

% For any local file, say "hello.tex" you want to link to please

% Theorem environments.
%
\theoremstyle{plain}
\newtheorem{theorem}[subsection]{Theorem}
\newtheorem{proposition}[subsection]{Proposition}
\newtheorem{lemma}[subsection]{Lemma}

\theoremstyle{definition}
\newtheorem{definition}[subsection]{Definition}
\newtheorem{example}[subsection]{Example}
\newtheorem{exercise}[subsection]{Exercise}
\newtheorem{situation}[subsection]{Situation}

\theoremstyle{remark}
\newtheorem{remark}[subsection]{Remark}
\newtheorem{remarks}[subsection]{Remarks}

\numberwithin{equation}{subsection}

% Macros
%
\def\lim{\mathop{\mathrm{lim}}\nolimits}
\def\colim{\mathop{\mathrm{colim}}\nolimits}
\def\Spec{\mathop{\mathrm{Spec}}}
\def\Hom{\mathop{\mathrm{Hom}}\nolimits}
\def\Ext{\mathop{\mathrm{Ext}}\nolimits}
\def\SheafHom{\mathop{\mathcal{H}\!\mathit{om}}\nolimits}
\def\SheafExt{\mathop{\mathcal{E}\!\mathit{xt}}\nolimits}
\def\Sch{\mathit{Sch}}
\def\Mor{\mathop{\mathrm{Mor}}\nolimits}
\def\Ob{\mathop{\mathrm{Ob}}\nolimits}
\def\Sh{\mathop{\mathit{Sh}}\nolimits}
\def\NL{\mathop{N\!L}\nolimits}
\def\CH{\mathop{\mathrm{CH}}\nolimits}
\def\proetale{{pro\text{-}\acute{e}tale}}
\def\etale{{\acute{e}tale}}
\def\QCoh{\mathit{QCoh}}
\def\Ker{\mathop{\mathrm{Ker}}}
\def\Im{\mathop{\mathrm{Im}}}
\def\Coker{\mathop{\mathrm{Coker}}}
\def\Coim{\mathop{\mathrm{Coim}}}

% Boxtimes
%
\DeclareMathSymbol{\boxtimes}{\mathbin}{AMSa}{"02}

%
% Macros for moduli stacks/spaces
%
\def\QCohstack{\mathcal{QC}\!\mathit{oh}}
\def\Cohstack{\mathcal{C}\!\mathit{oh}}
\def\Spacesstack{\mathcal{S}\!\mathit{paces}}
\def\Quotfunctor{\mathrm{Quot}}
\def\Hilbfunctor{\mathrm{Hilb}}
\def\Curvesstack{\mathcal{C}\!\mathit{urves}}
\def\Polarizedstack{\mathcal{P}\!\mathit{olarized}}
\def\Complexesstack{\mathcal{C}\!\mathit{omplexes}}
% \Pic is the operator that assigns to X its picard group, usage \Pic(X)
% \Picardstack_{X/B} denotes the Picard stack of X over B
% \Picardfunctor_{X/B} denotes the Picard functor of X over B
\def\Pic{\mathop{\mathrm{Pic}}\nolimits}
\def\Picardstack{\mathcal{P}\!\mathit{ic}}
\def\Picardfunctor{\mathrm{Pic}}
\def\Deformationcategory{\mathcal{D}\!\mathit{ef}}

\setlength{\emergencystretch}{3em}
\begin{document}
\title[Stacks Project, Tag 0BSR]{477 --- Étale-local product splitting of integral algebras with finite fibres and finite separable residue extensions}
\maketitle
\noindent Copyright (C) 2005--2025 Johan de Jong. Stacks Project authors. Source: \href{https://stacks.math.columbia.edu/tag/0BSR}{Tag 0BSR}.

\noindent Editorial difficulty note (not author text). Difficulty H1: The proof first captures every separable residue-field generator in a finite subalgebra while preserving distinct fibre primes. Étale splitting of that finite subalgebra must then be transferred to the original integral algebra, removing the unwanted summand without losing the purely inseparable residue-field condition..

\noindent Editorial provenance note (not author text). History: excerpt from revision \texttt{a0444\allowbreak 6e57e\allowbreak c1fbc\allowbreak 252a8\allowbreak 71afc\allowbreak ec775\allowbreak 2fb28\allowbreak 07b14}, more-morphisms.tex, lines 12210--12296. Reference presentation was adapted; original-author context and core support blocks were retained or added. The mathematical wording is unchanged. Licensed under GNU FDL 1.2 or later; no invariant sections or cover texts. See ../licenses/COPYING. Context blocks reproduce author definitions and supporting results; they are not additional counted problems.

\begin{lemma}
\label{lemma-etale-makes-integral-split}
Let $R \to S$ be an integral ring map. Let $\mathfrak p \subset R$ be a prime
ideal. Assume
\begin{enumerate}
\item there are finitely many primes $\mathfrak q_1, \ldots, \mathfrak q_n$
lying over $\mathfrak p$, and
\item for each $i$ the maximal separable subextension
$\kappa(\mathfrak q)/\kappa(\mathfrak q_i)_{sep}/\kappa(\mathfrak p)$
(Fields, Lemma \href{https://stacks.math.columbia.edu/tag/030K}{lemma separable first (Tag 030K)})
is finite over $\kappa(\mathfrak p)$.
\end{enumerate}
Then there exists an \'etale ring map $R \to R'$ and a prime
$\mathfrak p'$ lying over $\mathfrak p$ such that
$$
S \otimes_R R' = A_1 \times \ldots \times A_m
$$
with $R' \to A_j$ integral having a unique prime $\mathfrak r_j$
over $\mathfrak p'$ such that $\kappa(\mathfrak r_j)/\kappa(\mathfrak p')$
is purely inseparable.
\end{lemma}

\begin{proof}[First proof]
This proof uses
Algebra, Lemma \href{https://stacks.math.columbia.edu/tag/00UL}{lemma etale makes quasi finite finite variant (Tag 00UL)}.
Namely, choose a generator $\theta_i \in \kappa(\mathfrak q_i)_{sep}$
of this field over $\kappa(\mathfrak p)$
(Fields, Lemma \href{https://stacks.math.columbia.edu/tag/030N}{lemma primitive element (Tag 030N)}).
The spectrum of the fibre ring $S \otimes_R \kappa(\mathfrak p)$
is finite discrete with points corresponding to
$\mathfrak q_1, \ldots, \mathfrak q_n$.
By the Chinese remainder theorem
(Algebra, Lemma \href{https://stacks.math.columbia.edu/tag/00DT}{lemma chinese remainder (Tag 00DT)})
we see that $S \otimes_R \kappa(\mathfrak p) \to \prod \kappa(\mathfrak q_i)$
is surjective. Hence after replacing $R$ by $R_g$ for some
$g \in R$, $g \not \in \mathfrak p$ we may assume that
$(0, \ldots, 0, \theta_i, 0, \ldots, 0) \in \prod \kappa(\mathfrak q_i)$
is the image of some $x_i \in S$.
Let $S' \subset S$ be the $R$-subalgebra generated by our $x_i$.
Since $\Spec(S) \to \Spec(S')$ is surjective
(Algebra, Lemma \href{https://stacks.math.columbia.edu/tag/00GQ}{lemma integral overring surjective (Tag 00GQ)})
we conclude that
$\mathfrak q_i' = S' \cap \mathfrak q_i$ are the primes of
$S'$ over $\mathfrak p$. By our choice of $x_i$ we conclude
these primes are distinct that and
$\kappa(\mathfrak q'_i)_{sep} = \kappa(\mathfrak q_i)_{sep}$.
In particular the field extensions
$\kappa(\mathfrak q_i)/\kappa(\mathfrak q'_i)$ are purely
inseparable.
Since $R \to S'$ is finite we may apply
Algebra, Lemma \href{https://stacks.math.columbia.edu/tag/00UL}{lemma etale makes quasi finite finite variant (Tag 00UL)}.
and we get $R \to R'$ and $\mathfrak p'$ and a decomposition
$$
S' \otimes_R R' = A'_1 \times \ldots \times A'_m \times B'
$$
with $R' \to A'_j$ integral having a unique prime
$\mathfrak r'_j$ over $\mathfrak p'$ such that
$\kappa(\mathfrak r'_j)/\kappa(\mathfrak p')$
is purely inseparable and such that $B'$ does not have a prime
lying over $\mathfrak p'$. Since $R' \to B'$ is finite
(as $R \to S'$ is finite) we can
after localizing $R'$ at some $g' \in R'$, $g' \not \in \mathfrak p'$
assume that $B' = 0$. Via the map $S' \otimes_R R' \to S \otimes_R R'$
we get the corresponding decomposition for $S$.
\end{proof}

\begin{proof}[Second proof]
This proof uses strict henselization. First, assume $R$ is strictly
henselization with maximal ideal $\mathfrak p$. Then $S/\mathfrak p S$
has finitely many primes corresponding to
$\mathfrak q_1, \ldots, \mathfrak q_n$, each maximal,
each with purely inseparable residue field over $\kappa(\mathfrak p)$.
Hence $S/\mathfrak p S$ is equal to $\prod (S/\mathfrak p S)_{\mathfrak p_i}$.
By More on Algebra, Lemma \href{https://stacks.math.columbia.edu/tag/09XI}{lemma characterize henselian pair (Tag 09XI)}
we can lift this product decomposition to a product composition of $S$
as in the statement.

\medskip\noindent
In the general case, let $R^{sh}$ be the strict henselization of
$R_\mathfrak p$. Then we can apply the result of the first paragraph
to $R^{sh} \to S \otimes_R R^{sh}$. Consider the $m$ mutually orthogonal
idempotents in $S \otimes_R R^{sh}$ corresponding to the product
decomposition. Since $R^{sh}$ is a filtered colimit of \'etale
ring maps $(R, \mathfrak p) \to (R', \mathfrak p')$ by
Algebra, Lemma \href{https://stacks.math.columbia.edu/tag/04GW}{lemma strict henselization different (Tag 04GW)}
we see that these idempotents descend to some $R'$ as desired.
\end{proof}
\end{document}
